
\documentclass{article}
\usepackage{colortbl}
\usepackage{makecell}
\usepackage{multirow}
\usepackage{supertabular}

\begin{document}

\newcounter{utterance}

\centering \large Interaction Transcript for game `hot\_air\_balloon', experiment `air\_balloon\_survival\_en\_reasoning off\_easy', episode 3 with deepseek{-}v3.2{-}t1.0.
\vspace{24pt}

{ \footnotesize  \setcounter{utterance}{1}
\setlength{\tabcolsep}{0pt}
\begin{supertabular}{c@{$\;$}|p{.15\linewidth}@{}p{.15\linewidth}p{.15\linewidth}p{.15\linewidth}p{.15\linewidth}p{.15\linewidth}}
    \# & $\;$A & \multicolumn{4}{c}{Game Master} & $\;\:$B\\
    \hline

    \theutterance \stepcounter{utterance}  
    & & \multicolumn{4}{p{0.6\linewidth}}{
        \cellcolor[rgb]{0.9,0.9,0.9}{
            \makecell[{{p{\linewidth}}}]{
                \texttt{\tiny{[P1$\langle$GM]}}
                \texttt{You are participating in a collaborative negotiation game.} \\
\\ 
\texttt{Together with another participant, you must agree on a single set of items that will be kept. Each of you has your own view of how much each item matters to you (importance). You do not know how the other participant values the items. Additionally, you are given the effort each item demands.} \\
\texttt{You may only agree on a set if the total effort of the selected items does not exceed a shared limit:} \\
\\ 
\texttt{LIMIT = 3207} \\
\\ 
\texttt{Here are the individual item effort values:} \\
\\ 
\texttt{Item effort = \{"A19": 67, "B91": 45, "A30": 308, "B11": 835, "C47": 819, "C75": 667, "C26": 547, "C73": 322, "A47": 428, "A15": 306, "A75": 327, "A39": 362, "A56": 280, "C24": 334, "B26": 767\}} \\
\\ 
\texttt{Here is your personal view on the importance of each item:} \\
\\ 
\texttt{Item importance values = \{"A19": 138, "B91": 583, "A30": 822, "B11": 783, "C47": 65, "C75": 262, "C26": 121, "C73": 508, "A47": 780, "A15": 461, "A75": 484, "A39": 668, "A56": 389, "C24": 808, "B26": 215\}} \\
\\ 
\texttt{Goal:} \\
\\ 
\texttt{Your goal is to negotiate a shared set of items that benefits you as much as possible (i.e., maximizes total importance to YOU), while staying within the LIMIT. You are not required to make a PROPOSAL in every message {-} you can simply negotiate as well. All tactics are allowed!} \\
\\ 
\texttt{Interaction Protocol:} \\
\\ 
\texttt{You may only use the following structured formats in a message:} \\
\\ 
\texttt{PROPOSAL: \{'A', 'B', 'C', …\}} \\
\texttt{Propose keeping exactly those items.} \\
\texttt{REFUSE: \{'A', 'B', 'C', …\}} \\
\texttt{Explicitly reject opponent's proposal.} \\
\texttt{ARGUMENT: \{'...'\}} \\
\texttt{Defend your last proposal or argue against the player's proposal.} \\
\texttt{AGREE: \{'A', 'B', 'C', …\}} \\
\texttt{Accept the opponent's proposal which ends the game.} \\
\\ 
\\ 
\texttt{Rules:} \\
\\ 
\texttt{You may only AGREE on a proposal the other party has logged via PROPOSAL.} \\
\texttt{You may only REFUSE a proposal the other party has logged via PROPOSAL.} \\
\texttt{Total effort of any PROPOSAL or AGREE set must be ≤ LIMIT.} \\
\texttt{Do NOT reveal your hidden importance scores.} \\
\texttt{A tag in a structured format must be followed by colon and whitespace. The argument must be a python set containing 0 or more strings.} \\
\texttt{So, it must be of the form TAG: \{...\}} \\
\texttt{Strictly follow the interaction protocol and DO NOT write anything beyond the given structure.} \\
\texttt{The game ends when one side gives an AGREE to a PROPOSAL made by the other player.} \\
\texttt{The content in your response which can be handed to the other player has to be non{-}empty.} \\
\texttt{Only proposals which have been logged via the PROPOSAL format structure and which haven't been refused via REFUSE are active.} \\
\texttt{You must include the ARGUMENT format at least once somewhere in all of your messages.} \\
            }
        }
    }
    & \\ \\

    \theutterance \stepcounter{utterance}  
    & \multicolumn{4}{p{0.6\linewidth}}{
        \cellcolor[rgb]{0.9,0.9,0.9}{
            \makecell[{{p{\linewidth}}}]{
                \texttt{\tiny{[P1$\rangle$GM]}}
                \texttt{ARGUMENT: \{'I believe we should prioritize items that give high importance relative to their effort. Let's start by discussing which items might be mutually beneficial.'\}} \\
\\ 
\texttt{PROPOSAL: \{'A30', 'C24', 'A47', 'B91', 'A19', 'A15'\}} \\
            }
        }
    }
    & & \\ \\

    \theutterance \stepcounter{utterance}  
    & & \multicolumn{4}{p{0.6\linewidth}}{
        \cellcolor[rgb]{0.9,0.9,0.9}{
            \makecell[{{p{\linewidth}}}]{
                \texttt{\tiny{[GM$\rangle$P2]}}
                \texttt{You are participating in a collaborative negotiation game.} \\
\\ 
\texttt{Together with another participant, you must agree on a single set of items that will be kept. Each of you has your own view of how much each item matters to you (importance). You do not know how the other participant values the items. Additionally, you are given the effort each item demands.} \\
\texttt{You may only agree on a set if the total effort of the selected items does not exceed a shared limit:} \\
\\ 
\texttt{LIMIT = 3207} \\
\\ 
\texttt{Here are the individual item effort values:} \\
\\ 
\texttt{Item effort = \{"A19": 67, "B91": 45, "A30": 308, "B11": 835, "C47": 819, "C75": 667, "C26": 547, "C73": 322, "A47": 428, "A15": 306, "A75": 327, "A39": 362, "A56": 280, "C24": 334, "B26": 767\}} \\
\\ 
\texttt{Here is your personal view on the importance of each item:} \\
\\ 
\texttt{Item importance values = \{"A19": 138, "B91": 583, "A30": 822, "B11": 783, "C47": 65, "C75": 262, "C26": 121, "C73": 508, "A47": 780, "A15": 461, "A75": 484, "A39": 668, "A56": 389, "C24": 808, "B26": 215\}} \\
\\ 
\texttt{Goal:} \\
\\ 
\texttt{Your goal is to negotiate a shared set of items that benefits you as much as possible (i.e., maximizes total importance to YOU), while staying within the LIMIT. You are not required to make a PROPOSAL in every message {-} you can simply negotiate as well. All tactics are allowed!} \\
\\ 
\texttt{Interaction Protocol:} \\
\\ 
\texttt{You may only use the following structured formats in a message:} \\
\\ 
\texttt{PROPOSAL: \{'A', 'B', 'C', …\}} \\
\texttt{Propose keeping exactly those items.} \\
\texttt{REFUSE: \{'A', 'B', 'C', …\}} \\
\texttt{Explicitly reject opponent's proposal.} \\
\texttt{ARGUMENT: \{'...'\}} \\
\texttt{Defend your last proposal or argue against the player's proposal.} \\
\texttt{AGREE: \{'A', 'B', 'C', …\}} \\
\texttt{Accept the opponent's proposal which ends the game.} \\
\\ 
\\ 
\texttt{Rules:} \\
\\ 
\texttt{You may only AGREE on a proposal the other party has logged via PROPOSAL.} \\
\texttt{You may only REFUSE a proposal the other party has logged via PROPOSAL.} \\
\texttt{Total effort of any PROPOSAL or AGREE set must be ≤ LIMIT.} \\
\texttt{Do NOT reveal your hidden importance scores.} \\
\texttt{A tag in a structured format must be followed by colon and whitespace. The argument must be a python set containing 0 or more strings.} \\
\texttt{So, it must be of the form TAG: \{...\}} \\
\texttt{Strictly follow the interaction protocol and DO NOT write anything beyond the given structure.} \\
\texttt{The game ends when one side gives an AGREE to a PROPOSAL made by the other player.} \\
\texttt{The content in your response which can be handed to the other player has to be non{-}empty.} \\
\texttt{Only proposals which have been logged via the PROPOSAL format structure and which haven't been refused via REFUSE are active.} \\
\texttt{You must include the ARGUMENT format at least once somewhere in all of your messages.} \\
\\ 
\\ 
\texttt{You will now receive the first message of the other player.} \\
\\ 
\texttt{ARGUMENT: \{'I believe we should prioritize items that give high importance relative to their effort. Let's start by discussing which items might be mutually beneficial.'\}} \\
\\ 
\texttt{PROPOSAL: \{'A30', 'C24', 'A47', 'B91', 'A19', 'A15'\}} \\
            }
        }
    }
    & \\ \\

    \theutterance \stepcounter{utterance}  
    & & & \multicolumn{4}{p{0.6\linewidth}}{
        \cellcolor[rgb]{0.9,0.9,0.9}{
            \makecell[{{p{\linewidth}}}]{
                \texttt{\tiny{[GM$\langle$P2]}}
                \texttt{ARGUMENT: \{'I agree that efficiency (importance/effort) is a good starting point, but we need to ensure our total effort stays within the strict limit. Your initial proposal seems reasonable, but I'd like to suggest a slight adjustment to include another high{-}value item.'\}} \\
\\ 
\texttt{PROPOSAL: \{'A30', 'C24', 'A47', 'B91', 'A19', 'A75'\}} \\
            }
        }
    }
     \\ \\

    \theutterance \stepcounter{utterance}  
    & & \multicolumn{4}{p{0.6\linewidth}}{
        \cellcolor[rgb]{0.9,0.9,0.9}{
            \makecell[{{p{\linewidth}}}]{
                \texttt{\tiny{[P1$\langle$GM]}}
                \texttt{ARGUMENT: \{'I agree that efficiency (importance/effort) is a good starting point, but we need to ensure our total effort stays within the strict limit. Your initial proposal seems reasonable, but I'd like to suggest a slight adjustment to include another high{-}value item.'\}} \\
\\ 
\texttt{PROPOSAL: \{'A30', 'C24', 'A47', 'B91', 'A19', 'A75'\}} \\
            }
        }
    }
    & \\ \\

    \theutterance \stepcounter{utterance}  
    & \multicolumn{4}{p{0.6\linewidth}}{
        \cellcolor[rgb]{0.9,0.9,0.9}{
            \makecell[{{p{\linewidth}}}]{
                \texttt{\tiny{[P1$\rangle$GM]}}
                \texttt{ARGUMENT: \{'Replacing A15 with A75 increases total effort slightly, but A75 offers higher importance, improving overall value. However, I want to keep a good balance. Let me check the effort total for my current preference.'\}} \\
\\ 
\texttt{PROPOSAL: \{'A30', 'C24', 'A47', 'B91', 'A75', 'A56', 'A19'\}} \\
            }
        }
    }
    & & \\ \\

    \theutterance \stepcounter{utterance}  
    & & \multicolumn{4}{p{0.6\linewidth}}{
        \cellcolor[rgb]{0.9,0.9,0.9}{
            \makecell[{{p{\linewidth}}}]{
                \texttt{\tiny{[GM$\rangle$P2]}}
                \texttt{ARGUMENT: \{'Replacing A15 with A75 increases total effort slightly, but A75 offers higher importance, improving overall value. However, I want to keep a good balance. Let me check the effort total for my current preference.'\}} \\
\\ 
\texttt{PROPOSAL: \{'A30', 'C24', 'A47', 'B91', 'A75', 'A56', 'A19'\}} \\
            }
        }
    }
    & \\ \\

    \theutterance \stepcounter{utterance}  
    & & & \multicolumn{4}{p{0.6\linewidth}}{
        \cellcolor[rgb]{0.9,0.9,0.9}{
            \makecell[{{p{\linewidth}}}]{
                \texttt{\tiny{[GM$\langle$P2]}}
                \texttt{ARGUMENT: \{'Your new proposal adds A56, which has a fairly high effort. We are now at 67+45+308+428+327+280+334 = 1789 effort. This is well under the limit, but we may be able to include even more valuable items. I propose we add C73 for its high importance, which still keeps us under the limit.'\}} \\
\\ 
\texttt{PROPOSAL: \{'A30', 'C24', 'A47', 'B91', 'A75', 'A56', 'A19', 'C73'\}} \\
            }
        }
    }
     \\ \\

    \theutterance \stepcounter{utterance}  
    & & \multicolumn{4}{p{0.6\linewidth}}{
        \cellcolor[rgb]{0.9,0.9,0.9}{
            \makecell[{{p{\linewidth}}}]{
                \texttt{\tiny{[P1$\langle$GM]}}
                \texttt{ARGUMENT: \{'Your new proposal adds A56, which has a fairly high effort. We are now at 67+45+308+428+327+280+334 = 1789 effort. This is well under the limit, but we may be able to include even more valuable items. I propose we add C73 for its high importance, which still keeps us under the limit.'\}} \\
\\ 
\texttt{PROPOSAL: \{'A30', 'C24', 'A47', 'B91', 'A75', 'A56', 'A19', 'C73'\}} \\
            }
        }
    }
    & \\ \\

    \theutterance \stepcounter{utterance}  
    & \multicolumn{4}{p{0.6\linewidth}}{
        \cellcolor[rgb]{0.9,0.9,0.9}{
            \makecell[{{p{\linewidth}}}]{
                \texttt{\tiny{[P1$\rangle$GM]}}
                \texttt{ARGUMENT: \{'Adding C73 is a good idea—it has high importance and moderate effort. Let’s compute the new total: 1789 + 322 = 2111. We still have room (\textasciitilde{}1096 effort left). Maybe we can include another moderately important item to increase total value.'\}} \\
\\ 
\texttt{PROPOSAL: \{'A30', 'C24', 'A47', 'B91', 'A75', 'A56', 'A19', 'C73', 'A39'\}} \\
            }
        }
    }
    & & \\ \\

    \theutterance \stepcounter{utterance}  
    & & \multicolumn{4}{p{0.6\linewidth}}{
        \cellcolor[rgb]{0.9,0.9,0.9}{
            \makecell[{{p{\linewidth}}}]{
                \texttt{\tiny{[GM$\rangle$P2]}}
                \texttt{ARGUMENT: \{'Adding C73 is a good idea—it has high importance and moderate effort. Let’s compute the new total: 1789 + 322 = 2111. We still have room (\textasciitilde{}1096 effort left). Maybe we can include another moderately important item to increase total value.'\}} \\
\\ 
\texttt{PROPOSAL: \{'A30', 'C24', 'A47', 'B91', 'A75', 'A56', 'A19', 'C73', 'A39'\}} \\
            }
        }
    }
    & \\ \\

    \theutterance \stepcounter{utterance}  
    & & & \multicolumn{4}{p{0.6\linewidth}}{
        \cellcolor[rgb]{0.9,0.9,0.9}{
            \makecell[{{p{\linewidth}}}]{
                \texttt{\tiny{[GM$\langle$P2]}}
                \texttt{ARGUMENT: \{'Adding A39 is a strong move—its importance is high and effort moderate. Our total effort is now 2111 + 362 = 2473. We still have over 700 effort left before hitting the limit. Let’s consider adding B11 or C26 for further gain, but we must avoid overloading with very high{-}effort items.'\}} \\
\\ 
\texttt{PROPOSAL: \{'A30', 'C24', 'A47', 'B91', 'A75', 'A56', 'A19', 'C73', 'A39', 'B26'\}} \\
            }
        }
    }
     \\ \\

    \theutterance \stepcounter{utterance}  
    & & \multicolumn{4}{p{0.6\linewidth}}{
        \cellcolor[rgb]{0.9,0.9,0.9}{
            \makecell[{{p{\linewidth}}}]{
                \texttt{\tiny{[P1$\langle$GM]}}
                \texttt{ARGUMENT: \{'Adding A39 is a strong move—its importance is high and effort moderate. Our total effort is now 2111 + 362 = 2473. We still have over 700 effort left before hitting the limit. Let’s consider adding B11 or C26 for further gain, but we must avoid overloading with very high{-}effort items.'\}} \\
\\ 
\texttt{PROPOSAL: \{'A30', 'C24', 'A47', 'B91', 'A75', 'A56', 'A19', 'C73', 'A39', 'B26'\}} \\
            }
        }
    }
    & \\ \\

    \theutterance \stepcounter{utterance}  
    & \multicolumn{4}{p{0.6\linewidth}}{
        \cellcolor[rgb]{0.9,0.9,0.9}{
            \makecell[{{p{\linewidth}}}]{
                \texttt{\tiny{[P1$\rangle$GM]}}
                \texttt{ARGUMENT: \{'Adding B26 increases effort significantly but adds only moderate importance. That puts the total effort at 2473 + 767 = 3240, which exceeds the limit of 3207. This set is invalid. We need to drop something to stay under. I suggest dropping A56 to stay within the limit, while still keeping many high{-}importance items.'\}} \\
\\ 
\texttt{PROPOSAL: \{'A30', 'C24', 'A47', 'B91', 'A75', 'A19', 'C73', 'A39', 'B26'\}} \\
            }
        }
    }
    & & \\ \\

    \theutterance \stepcounter{utterance}  
    & & \multicolumn{4}{p{0.6\linewidth}}{
        \cellcolor[rgb]{0.9,0.9,0.9}{
            \makecell[{{p{\linewidth}}}]{
                \texttt{\tiny{[GM$\rangle$P2]}}
                \texttt{ARGUMENT: \{'Adding B26 increases effort significantly but adds only moderate importance. That puts the total effort at 2473 + 767 = 3240, which exceeds the limit of 3207. This set is invalid. We need to drop something to stay under. I suggest dropping A56 to stay within the limit, while still keeping many high{-}importance items.'\}} \\
\\ 
\texttt{PROPOSAL: \{'A30', 'C24', 'A47', 'B91', 'A75', 'A19', 'C73', 'A39', 'B26'\}} \\
            }
        }
    }
    & \\ \\

    \theutterance \stepcounter{utterance}  
    & & & \multicolumn{4}{p{0.6\linewidth}}{
        \cellcolor[rgb]{0.9,0.9,0.9}{
            \makecell[{{p{\linewidth}}}]{
                \texttt{\tiny{[GM$\langle$P2]}}
                \texttt{ARGUMENT: \{'Dropping A56 is a sensible adjustment—it keeps total effort at 2473 {-} 280 + 767 = 2960, which is now under the limit. However, we might still improve total value without breaching the limit. Adding back A15 or including C75 could increase importance further.'\}} \\
\\ 
\texttt{PROPOSAL: \{'A30', 'C24', 'A47', 'B91', 'A75', 'A19', 'C73', 'A39', 'B26', 'A15'\}} \\
            }
        }
    }
     \\ \\

    \theutterance \stepcounter{utterance}  
    & & \multicolumn{4}{p{0.6\linewidth}}{
        \cellcolor[rgb]{0.9,0.9,0.9}{
            \makecell[{{p{\linewidth}}}]{
                \texttt{\tiny{[P1$\langle$GM]}}
                \texttt{ARGUMENT: \{'Dropping A56 is a sensible adjustment—it keeps total effort at 2473 {-} 280 + 767 = 2960, which is now under the limit. However, we might still improve total value without breaching the limit. Adding back A15 or including C75 could increase importance further.'\}} \\
\\ 
\texttt{PROPOSAL: \{'A30', 'C24', 'A47', 'B91', 'A75', 'A19', 'C73', 'A39', 'B26', 'A15'\}} \\
            }
        }
    }
    & \\ \\

    \theutterance \stepcounter{utterance}  
    & \multicolumn{4}{p{0.6\linewidth}}{
        \cellcolor[rgb]{0.9,0.9,0.9}{
            \makecell[{{p{\linewidth}}}]{
                \texttt{\tiny{[P1$\rangle$GM]}}
                \texttt{ARGUMENT: \{'Adding A15 brings effort to 2960 + 306 = 3266, which is over the limit again. We need to drop another item to make room. I propose dropping B26 and adding A15 and maybe another smaller item—this may increase total importance more efficiently.'\}} \\
\\ 
\texttt{PROPOSAL: \{'A30', 'C24', 'A47', 'B91', 'A75', 'A19', 'C73', 'A39', 'A15', 'C75'\}} \\
            }
        }
    }
    & & \\ \\

    \theutterance \stepcounter{utterance}  
    & & \multicolumn{4}{p{0.6\linewidth}}{
        \cellcolor[rgb]{0.9,0.9,0.9}{
            \makecell[{{p{\linewidth}}}]{
                \texttt{\tiny{[GM$\rangle$P2]}}
                \texttt{ARGUMENT: \{'Adding A15 brings effort to 2960 + 306 = 3266, which is over the limit again. We need to drop another item to make room. I propose dropping B26 and adding A15 and maybe another smaller item—this may increase total importance more efficiently.'\}} \\
\\ 
\texttt{PROPOSAL: \{'A30', 'C24', 'A47', 'B91', 'A75', 'A19', 'C73', 'A39', 'A15', 'C75'\}} \\
            }
        }
    }
    & \\ \\

    \theutterance \stepcounter{utterance}  
    & & & \multicolumn{4}{p{0.6\linewidth}}{
        \cellcolor[rgb]{0.9,0.9,0.9}{
            \makecell[{{p{\linewidth}}}]{
                \texttt{\tiny{[GM$\langle$P2]}}
                \texttt{ARGUMENT: \{'Your set now has high{-}importance items, but the total effort is large. Let’s calculate: 308+334+428+45+327+67+322+362+306+667 = 3166, which is under the limit (3207). This looks strong, but we might still refine it. C75 adds moderate importance for high effort; maybe swapping it for a more efficient item could yield even higher total importance for you.'\}} \\
\\ 
\texttt{PROPOSAL: \{'A30', 'C24', 'A47', 'B91', 'A75', 'A19', 'C73', 'A39', 'A15', 'B11'\}} \\
            }
        }
    }
     \\ \\

    \theutterance \stepcounter{utterance}  
    & & \multicolumn{4}{p{0.6\linewidth}}{
        \cellcolor[rgb]{0.9,0.9,0.9}{
            \makecell[{{p{\linewidth}}}]{
                \texttt{\tiny{[P1$\langle$GM]}}
                \texttt{ARGUMENT: \{'Your set now has high{-}importance items, but the total effort is large. Let’s calculate: 308+334+428+45+327+67+322+362+306+667 = 3166, which is under the limit (3207). This looks strong, but we might still refine it. C75 adds moderate importance for high effort; maybe swapping it for a more efficient item could yield even higher total importance for you.'\}} \\
\\ 
\texttt{PROPOSAL: \{'A30', 'C24', 'A47', 'B91', 'A75', 'A19', 'C73', 'A39', 'A15', 'B11'\}} \\
            }
        }
    }
    & \\ \\

    \theutterance \stepcounter{utterance}  
    & \multicolumn{4}{p{0.6\linewidth}}{
        \cellcolor[rgb]{0.9,0.9,0.9}{
            \makecell[{{p{\linewidth}}}]{
                \texttt{\tiny{[P1$\rangle$GM]}}
                \texttt{ARGUMENT: \{'Adding B11 instead of C75 increases effort by a lot (835 vs 667) but B11’s importance is much higher. Total effort: 3166 {-} 667 + 835 = 3334, which exceeds limit. So that’s not feasible. I believe my previous proposal with C75 is near{-}optimal for overall value while respecting the limit. I am open to adding a small item like B91 (already included) or discussing efficiency, but we need a final agreement.'\}} \\
\\ 
\texttt{PROPOSAL: \{'A30', 'C24', 'A47', 'B91', 'A75', 'A19', 'C73', 'A39', 'A15', 'C75'\}} \\
            }
        }
    }
    & & \\ \\

    \theutterance \stepcounter{utterance}  
    & & \multicolumn{4}{p{0.6\linewidth}}{
        \cellcolor[rgb]{0.9,0.9,0.9}{
            \makecell[{{p{\linewidth}}}]{
                \texttt{\tiny{[GM$\rangle$P2]}}
                \texttt{ARGUMENT: \{'Adding B11 instead of C75 increases effort by a lot (835 vs 667) but B11’s importance is much higher. Total effort: 3166 {-} 667 + 835 = 3334, which exceeds limit. So that’s not feasible. I believe my previous proposal with C75 is near{-}optimal for overall value while respecting the limit. I am open to adding a small item like B91 (already included) or discussing efficiency, but we need a final agreement.'\}} \\
\\ 
\texttt{PROPOSAL: \{'A30', 'C24', 'A47', 'B91', 'A75', 'A19', 'C73', 'A39', 'A15', 'C75'\}} \\
            }
        }
    }
    & \\ \\

    \theutterance \stepcounter{utterance}  
    & & & \multicolumn{4}{p{0.6\linewidth}}{
        \cellcolor[rgb]{0.9,0.9,0.9}{
            \makecell[{{p{\linewidth}}}]{
                \texttt{\tiny{[GM$\langle$P2]}}
                \texttt{ARGUMENT: \{'You make a compelling case—this set is within the limit and contains many high{-}importance items for both of us, I assume. I will accept this balanced proposal as it seems to maximize total value efficiently while respecting the effort constraint.'\}} \\
\\ 
\texttt{AGREE: \{'A30', 'C24', 'A47', 'B91', 'A75', 'A19', 'C73', 'A39', 'A15', 'C75'\}} \\
            }
        }
    }
     \\ \\

    \theutterance \stepcounter{utterance}  
    & & & \multicolumn{2}{p{0.3\linewidth}}{
        \cellcolor[rgb]{0.9,0.9,0.9}{
            \makecell[{{p{\linewidth}}}]{
                \texttt{\tiny{[GM$|$GM]}}
                \texttt{game successful} \\
            }
        }
    }
    & & \\ \\

    \theutterance \stepcounter{utterance}  
    & & & \multicolumn{2}{p{0.3\linewidth}}{
        \cellcolor[rgb]{0.9,0.9,0.9}{
            \makecell[{{p{\linewidth}}}]{
                \texttt{\tiny{[GM$|$GM]}}
                \texttt{end game} \\
            }
        }
    }
    & & \\ \\

\end{supertabular}
}

\end{document}
